\section{Finite polynomials and the cost of assigning a color}
\label{sec:polynomials}

We develop a finite-dimensional estimate for assigning colors to the sites of a
Hermitian matrix.  The polynomial associated with a partial coloring will
average over the remaining color choices and over independent signs.  We
use the sum of the positive roots after a negative diagonal shift as a cost.
Assigning a color restricts the average and makes the shift more negative at
the chosen site; we will show that some color increases this cost by at most
a fixed multiple of a quantity attached to that site.  The locality and uniform bounds proved
alongside this estimate will permit its use on measured graphs.

\subsection{Finite data and mixed determinantal bounds}

Fix a finite set \(I\), a zero-diagonal Hermitian matrix
\(T=(T_{uv})_{u,v\in I}\) with \(\|T\|\leq1\), an integer \(r\geq1\), and
\(t>0\).  Throughout the coloring argument we use
\[
                         \beta=80.
\]
A partial coloring is a function \(c:I\to\{0,1,\ldots,r\}\), where \(0\)
means unassigned.  Give each site a weight \(W_v>0\), and put
\[
 \tau_v=
 \begin{cases}
 t,&c(v)=0,\\
 \beta t,&c(v)>0,
 \end{cases}
 \qquad a_v=\tau_vW_v.
\]
At site \(v\), the allowed pairs of color and sign are
\[
 \mathcal A_v=
 \begin{cases}
 \{1,\ldots,r\}\times\{+1,-1\},&c(v)=0,\\
 \{c(v)\}\times\{+1,-1\},&c(v)>0.
 \end{cases}
\]
Choose a pair \(\omega_v\) uniformly from \(\mathcal A_v\), independently
over sites, and write \(k_v=|\mathcal A_v|\).  For
\(j=(i,\sigma)\in\{1,\ldots,r\}\times\{+1,-1\}\), set
\[
 U_j=\{v\in I:j\in\mathcal A_v\},\qquad
 B_j=\sigma D_W T D_W-\diag(a_v)_{v\in I},
 \qquad D_W=\diag(\sqrt{W_v})_{v\in I}.
\]
For a realization \(\omega\), let \(I_j(\omega)=\{v:\omega_v=j\}\).
Its outcome matrix is the direct sum of the principal matrices
\(B_j|_{I_j(\omega)}\), placed on these subsets of \(I\); its entries
between different subsets are zero.  Denote this matrix by \(B_\omega\).
With \(Z=\diag(z_v)_{v\in I}\), define
\begin{equation}
 P_I(Z)=\mathbb E_\omega\det(Z-B_\omega),\qquad
 P_I(z)=P_I(z,\ldots,z).
 \label{eq:poly-definition}
\end{equation}
We write \(P=P_I\) when \(I\) is fixed.  Restriction to a subset always
means principal restriction of all the finite data, with the same laws at
the remaining sites.  In particular \(P_{-v}=P_{I\setminus\{v\}}\), and
\(P_\varnothing=1\).  These conventions also apply to all polynomials
introduced below.

We shall use the following two bounds from mixed determinantal polynomial
theory.  We state their precise normalization because both signs of \(T\)
will enter our application.

\begin{lemma}[Mixed determinantal bounds]\label{lem:mixed-bounds}
For a \(k\)-tuple of Hermitian matrices \(A_1,\ldots,A_k\) on a common
nonempty finite set, let
\[
 \chi[A_1,\ldots,A_k](z)
 =\mathbb E\prod_{j=1}^k\det(zI-A_j|_{S_j}),
\]
where each site is independently assigned uniformly to one of the \(k\)
sets \(S_j\), and the determinant of an empty matrix is \(1\).
This polynomial is real-rooted.  If \(k\geq2\) and the \(A_j\) are
zero-diagonal contractions, then
\begin{equation}
             \lambda_{\max}\chi[A_1,\ldots,A_k]
                        <\frac{3\sqrt2}{\sqrt{k}}.
 \label{eq:poly-mixed-upper}
\end{equation}
For any zero-diagonal Hermitian matrix \(L\), without a norm restriction,
\begin{equation}
                \lambda_{\max}\chi[L,-L]\geq\frac{\|L\|}{2}.
 \label{eq:poly-mixed-lower}
\end{equation}
Consequently, on a nonempty finite set, if every site is unassigned, the unit-weight polynomial
without diagonal shifts has largest root \(<3/\sqrt r\).  If every site
is assigned, that unit-weight unshifted polynomial is
\[
                   \prod_{i=1}^r\chi[T|_{c^{-1}(i)},-T|_{c^{-1}(i)}],
\]
where factors on empty color classes are \(1\).
\end{lemma}

\begin{proof}
The normalization agrees with Definition~5 of \cite{RS17}.
Real-rootedness is part of the mixed determinantal theory and also follows
from Lemma~\ref{lem:polynomial-quotients} below.
Equation~\eqref{eq:poly-mixed-upper} is
\cite[Theorem~25]{RS17}, whose hypotheses are zero-diagonal Hermitian
contractions and \(k\geq2\).

By \cite[Proposition~16 and Theorem~15]{RS17}, for a tuple of
zero-diagonal Hermitian matrices, replacing one member by zero cannot
increase its largest root.
Expansion in principal minors gives
\[
                    \chi[L,0](z)=\det(zI-L/2):
\]
a fixed principal minor on \(S\) occurs with probability \(2^{-|S|}\).
Applying monotonicity while retaining either \(L\) or \(-L\) proves
Equation~\eqref{eq:poly-mixed-lower}.  The all-unassigned case uses the
tuple containing \(r\) copies each of \(T\) and \(-T\), so \(k=2r\).
For a fully assigned coloring, independence between the color classes
gives the stated factorization.
\end{proof}

\subsection{Stable quotients and rooted spectral measures}

The next lemma provides the interlacing needed to measure the effect of
one site.  Its half-plane assertion also allows non-Hermitian matrices:
this extension will later control analytic perturbations of the data.
For a matrix \(B\), write
\(\operatorname{Im}B=(B-B^*)/(2i)\).

\begin{lemma}[Polynomial quotients]\label{lem:polynomial-quotients}
For the finite data above,
\begin{equation}
 P_I(Z)=
 \left(\prod_{v\in I}\frac{1}{k_v!}
                  \partial_{z_v}^{\,k_v-1}\right)
            \prod_j\det\bigl((Z-B_j)|_{U_j}\bigr),
 \qquad \partial_{z_v}P_I=P_{-v}.
 \label{eq:poly-differential}
\end{equation}
More generally, use any matrices \(B_j\) on \(U_j\) in the same
independent-choice construction.  If
\(\operatorname{Im}B_j\leq hI\) for a real number \(h\), then \(P_I(Z)\)
and every deletion polynomial are nonzero whenever
\(\operatorname{Im}z_v>h\) at every remaining site.  In this region,
\begin{equation}
 \operatorname{Im}\frac{P_I(Z)}{P_{-v}(Z|_{I\setminus\{v\}})}
                         \geq\operatorname{Im}z_v-h.
 \label{eq:poly-half-plane}
\end{equation}

Return to the Hermitian data in Equation~\eqref{eq:poly-definition}.
For \(v\in I\), set
\[
 Q_v=P_{-v},\qquad F_v(z)=\frac{P_I(z)}{Q_v(z)},
 \qquad G_v(z)=\frac{Q_v(z)}{P_I(z)}.
\]
These rational functions are analytic in the upper half-plane, and
\begin{equation}
 F_v(z)=z+a_v-\sum_{\lambda}\frac{b_{v,\lambda}}{z-\lambda},
                      \qquad b_{v,\lambda}\geq0,
 \label{eq:poly-arrow-quotient}
\end{equation}
where \(\lambda\) ranges over the distinct roots of \(Q_v\).
Let \(\Lambda_v\) be the diagonal matrix listing those roots with
multiplicity.  There is a vector
\(g_v\), with squared coordinate \(b_{v,\lambda}\) on one copy of each
\(\lambda\) and zero on its other copies, such that
\begin{equation}
                H_v=\begin{pmatrix}-a_v&g_v^*\\ g_v&\Lambda_v\end{pmatrix}
 \label{eq:poly-arrow}
\end{equation}
has characteristic polynomial \(P_I\).  Its first-coordinate spectral
measure \(\rho_v\) is a probability measure satisfying
\begin{equation}
 G_v(z)=\int_{\mathbb R}\frac{1}{z-x}\,\rho_v(dx),\qquad
 \sum_{v\in I}\rho_v=\sum_{P_I(x)=0}\delta_x.
 \label{eq:poly-root-measures}
\end{equation}
The roots on the right are counted with multiplicity, and \(Q_v\)
interlaces \(P_I\).
\end{lemma}

\begin{proof}
Each determinant factor is affine in any of its diagonal variables.
When differentiating \(k_v-1\) times in \(z_v\), every surviving Leibniz
term leaves \(v\) in exactly one of its \(k_v\) available factors and
deletes it from the others.  The multiplicity \((k_v-1)!\), divided by
\(k_v!\), gives probability \(1/k_v\).  Processing all vertices gives
the outcome average, and differentiation once more in \(z_v\) deletes
that site.  This proves Equation~\eqref{eq:poly-differential}, including
for the more general matrices.

If \(\operatorname{Im}z_v>h\), each matrix \(Z-B_j\) on \(U_j\) has
strictly positive imaginary part and is invertible.  Its determinant is
therefore nonzero.  Differentiation preserves nonvanishing in the
specified half-plane by the Gauss--Lucas Theorem.  Here the degree cannot
drop at a specialization of the other variables: after any subset of
the vertex operations has been performed, the leading coefficient of
degree \(k_v\) in an unprocessed variable \(z_v\) is the corresponding
partially processed expression with \(v\) deleted.  Induct simultaneously
over the processed vertices and all principal restrictions.  That
leading coefficient is nonzero by the induction hypothesis, so every
required derivative has its indicated degree and is nonzero in the
half-plane.  This also proves the assertion for all deletion polynomials.

For fixed other coordinates in that region,
\(P_I/P_{-v}=z_v+\zeta_0\), by the derivative identity and
multiaffineness.  Its zero has imaginary part at most \(h\).
This gives Equation~\eqref{eq:poly-half-plane}.

For Hermitian data, the uniform-variable polynomials have real
coefficients; half-plane nonvanishing and conjugation therefore show
that all their roots are real.  Equation~\eqref{eq:poly-half-plane}
with \(h=0\) implies that \(F_v(z)-z\) has nonnegative imaginary part.
Its constant at infinity is \(a_v\), since every outcome has \(v\)
diagonal \(-a_v\).  A nonremovable real pole of a real rational function
with nonnegative imaginary part must be simple with nonpositive
residue: approach the pole from different angles in the upper
half-plane to exclude higher orders, and then approach vertically to
determine the sign.  This proves
Equation~\eqref{eq:poly-arrow-quotient}.

The Schur complement of \(z-\Lambda_v\) in \(z-H_v\) is \(F_v(z)\).
Thus \(\det(z-H_v)=Q_v(z)F_v(z)=P_I(z)\), and its first resolvent
entry is \(1/F_v(z)=G_v(z)\).  The spectral theorem proves the first
identity in Equation~\eqref{eq:poly-root-measures}, and principal-matrix
interlacing proves the interlacing assertion.  Finally, uniform
specialization of Equation~\eqref{eq:poly-differential} gives
\[
                      \sum_{v\in I}G_v(z)=\frac{P_I'(z)}{P_I(z)}.
\]
The right side is the Cauchy transform of the root-counting measure.
Uniqueness of partial fractions proves the second identity in
Equation~\eqref{eq:poly-root-measures}.
\end{proof}

We will compare two modifications at an unassigned site \(v\).
For \(i\in\{1,\ldots,r\}\), first condition its pair on color \(i\),
keeping both signs uniformly distributed and retaining its shift
\(\tau_v=t\).  Denote the resulting polynomial and quotient by
\(P^0_{v,i}\) and \(F^0_{v,i}=P^0_{v,i}/Q_v\).
Next increase that site's shift to \(\beta t\), as required by a
permanent assignment.  Denote these by \(P_{v,i}\) and \(F_{v,i}\).
All other site data remain fixed, and deletion of \(v\) gives the same
polynomial \(Q_v\) in each case.  Consequently
\begin{equation}
 \frac1r\sum_{i=1}^r F^0_{v,i}=F_v,\qquad
 F_{v,i}=F^0_{v,i}+(\beta-1)tW_v.
 \label{eq:poly-options}
\end{equation}
The arrow representation and rooted measures apply to these modified
data as well.  In particular the residues of \(F^0_{v,i}\) are
nonnegative and average to the residues of \(F_v\).

\subsection{Positive-root costs and locality}

For a monic real-rooted polynomial \(P\), define
\[
 s(P)=\sum_{P(x)=0}x_+,\qquad x_+=\max\{x,0\},
 \qquad s(1)=0.
\]
The roots in this sum are counted with multiplicity.
For the finite site data, define the rooted cost and the deletion
increment by
\begin{equation}
       m_v=\int x_+\,\rho_v(dx),\qquad
       d_v=s(P_I)-s(P_{-v}).
 \label{eq:poly-costs}
\end{equation}
Thus \(\sum_vm_v=s(P_I)\).
Superscripts \(0,i\) and \(i\) will indicate, respectively, the
conditioned and permanently assigned options at \(v\).  For example,
\(d_v^{\,i}=s(P_{v,i})-s(Q_v)\).

The initial shift \(t>3/\sqrt r\) gives zero cost before any assignment.
For a fully assigned block, excessive pinched norm forces positive cost.
The following finite statement makes these two roles precise.

\begin{lemma}[Initial cost and detection of excessive norm]
\label{lem:finite-cost-endpoints}
For the finite data of this section, the following hold.
\begin{enumerate}
\item If every site is unassigned and \(t>3/\sqrt r\), then
\(P_I(x)>0\) for every \(x\geq0\), and \(s(P_I)=0\). This includes
\(I=\varnothing\), with \(P_\varnothing=1\).
\item Suppose \(I\ne\varnothing\) and every site is assigned, and put
\[
 T_c=\sum_{i=1}^r 1_{\{c=i\}}T1_{\{c=i\}}.
\]
For every \(\alpha>0\), if \(\|T_c\|>2\beta t+\alpha\), then
\[
 \lambda_{\max}P_I\geq\frac\alpha2\min_{v\in I}W_v,
 \qquad
 s(P_I)\geq\frac\alpha2\min_{v\in I}W_v.
\]
\end{enumerate}
\end{lemma}

\begin{proof}
For \(J\subseteq I\), let \(P_J^{\mathrm{un}}(Z)\) be the expected
polynomial with the same pair laws, all weights equal to one, and all
diagonal shifts removed. Factoring the weight matrix from every outcome
determinant gives
\begin{equation}\label{eq:poly-remove-weights}
 P_J(x)=\left(\prod_{v\in J}W_v\right)
 P_J^{\mathrm{un}}\bigl((\tau_v+x/W_v)_{v\in J}\bigr).
\end{equation}
Both sides equal one for \(J=\varnothing\).
We also use the following consequence of multiaffineness and the deletion
identity in Equation~\eqref{eq:poly-differential}. For real coordinates
\(b=(b_v)_{v\in I}\) and increments \(h_v\geq0\),
\[
 P_J^{\mathrm{un}}((b+h)|_J)
 =\sum_{S\subseteq J}
 P_{J\setminus S}^{\mathrm{un}}(b|_{J\setminus S})
       \prod_{v\in S}h_v.
\]
Thus positivity of all principal evaluations at \(b\) is preserved
when any coordinates are increased.

If every site is unassigned and \(t>3/\sqrt r\),
Lemma~\ref{lem:mixed-bounds} puts \(t\) strictly above the largest root
of every nonempty \(P_J^{\mathrm{un}}(z)\). All principal evaluations
at the common coordinate \(t\) are therefore positive. For \(x\geq0\),
the coordinates \(t+x/W_v\) in Equation~\eqref{eq:poly-remove-weights}
are at least \(t\), so \(P_I(x)>0\). Real-rootedness then gives
\(s(P_I)=0\).

Now suppose every site is assigned and \(\|T_c\|>2\beta t+\alpha\).
The fully assigned factorization and the lower bound in
Lemma~\ref{lem:mixed-bounds} show that \(P_I^{\mathrm{un}}(z)\) has a
root greater than \(\beta t+\alpha/2\). Set
\(x_0=(\alpha/2)\min_{v\in I}W_v>0\). If the largest root of \(P_I\)
were smaller than \(x_0\), repeated interlacing would put \(x_0\)
strictly above every root of every nonempty principal restriction.
Thus \(P_J(x_0)>0\) for all \(J\subseteq I\). By
Equation~\eqref{eq:poly-remove-weights}, all the corresponding
unweighted principal evaluations are positive at
\[
 b_v=\beta t+x_0/W_v\leq\beta t+\alpha/2.
\]
Increasing these coordinates to any common value
\(q\geq\beta t+\alpha/2\) preserves that positivity. In particular
\(P_I^{\mathrm{un}}(q)>0\) for every such \(q\), contradicting the root
above \(\beta t+\alpha/2\). Hence the largest root of \(P_I\) is at
least \(x_0\), which also gives the stated lower bound for \(s(P_I)\).
\end{proof}

\begin{lemma}[Deletion costs and their tails]\label{lem:positive-cost}
For the finite Hermitian data above, let
\(n_P(x)=\#\{\lambda:P(\lambda)=0,\ \lambda>x\}\), with multiplicity.
Then, for every \(v\in I\),
\begin{align}
 0&\leq n_{P_I}(x)-n_{Q_v}(x)\leq1, \label{eq:poly-interlacing-count}\\
 d_v&=\int_0^\infty(n_{P_I}-n_{Q_v})(x)\,dx, \label{eq:poly-cost-integral}\\
 \int_B^\infty(n_{P_I}-n_{Q_v})(x)\,dx
 &\leq 2\int(x-B)_+\,\rho_v(dx)
              +(a_v+B)\rho_v((B,\infty)),\qquad B\geq0.
 \label{eq:poly-cost-tail}
\end{align}
In particular \(0\leq d_v\leq2m_v+a_v\).
All these conclusions hold for either modification at \(v\), with its
own shift and rooted measure.
\end{lemma}

\begin{proof}
Equation~\eqref{eq:poly-interlacing-count} follows from the principal
interlacing in Lemma~\ref{lem:polynomial-quotients}; integration gives
Equation~\eqref{eq:poly-cost-integral}.  To prove the tail estimate,
write \(H=H_v\), \(\Lambda=\Lambda_v\), \(a=a_v\), and
\(L=H-BI\).  Let \(E\) be the first-coordinate projection,
\(D=1-E\), and \(R=1_{(0,\infty)}(L)\).
The compression \(DRD\) is a positive contraction on the deletion
coordinates.  The variational formula for the positive trace therefore
gives
\[
 \Tr(\Lambda-BI)_+\geq\Tr\bigl((\Lambda-BI)DRD\bigr).
\]
Since \(L-DLD=EL+LE-ELE\), it follows that
\begin{align*}
 \Tr(H-BI)_+-\Tr(\Lambda-BI)_+
 &\leq\Tr\bigl((EL+LE-ELE)R\bigr)\\
 &=2\langle e_0,LR e_0\rangle+(a+B)\langle e_0,R e_0\rangle\\
 &=2\int(x-B)_+\,\rho_v(dx)+(a+B)\rho_v((B,\infty)).
\end{align*}
Here \(e_0\) is the first unit vector of the arrow matrix.
The left side is the tail integral in
Equation~\eqref{eq:poly-cost-tail}.  Taking \(B=0\) and using that
\(\rho_v\) is a probability measure gives the last assertion.
\end{proof}

The deletion increment is defined using all roots of a polynomial,
whereas \(m_v\) is attached to one site.  We next record the locality
which relates these finite quantities to graph neighborhoods.  A graph
supporting \(T\) means an undirected graph on \(I\), without loops,
such that \(T_{uv}=0\) whenever \(u\ne v\) are not adjacent.

\begin{lemma}[Rooted locality and bounds]\label{lem:rooted-locality}
Let a graph support the finite zero-diagonal Hermitian contraction \(T\).
The coefficient of \(z^{-1-k}\) in \(G_v(z)\) is
\(\int x^k\,\rho_v(dx)\) and is determined by the data on the radius-\(k\)
ball about \(v\).  The Laurent coefficients of \(F_v(z)\), and of the
two options \(F^0_{v,i}(z)\) and \(F_{v,i}(z)\) at an unassigned site,
likewise depend only on a radius bounded by their order.
Moreover,
\begin{equation}
 \int x^2\,\rho_v(dx)
 =a_v^2+\sum_{u\ne v}W_vW_u|T_{vu}|^2
                   \mathbb P(\omega_u=\omega_v),
 \qquad
 \sum_\lambda b_{v,\lambda}
 =\sum_{u\ne v}W_vW_u|T_{vu}|^2
                   \mathbb P(\omega_u=\omega_v).
 \label{eq:poly-second-moment}
\end{equation}
The same identities hold for either site option, using the corresponding
law of the pairs and diagonal shift.

If neighboring weights have ratios in \([K^{-1},K]\), \(K\geq1\), then
for all these data
\begin{equation}
 \int x^2\,\rho_v(dx)\leq\bigl((\beta t)^2+K\bigr)W_v^2,\quad
 m_v\leq C_0W_v,\quad d_v\leq(2C_0+\beta t)W_v,
 \quad C_0=\sqrt{(\beta t)^2+K}.
 \label{eq:poly-weight-bounds}
\end{equation}
If all weights are at most \(H_0\), every polynomial and every principal
restriction has all its roots in
\begin{equation}
                     [-L,L],\qquad L=(1+\beta t)H_0.
 \label{eq:poly-root-support}
\end{equation}
The rooted measures have the same support bound.
\end{lemma}

\begin{proof}
Put \(u_v=z_v^{-1}\) and consider the formal polynomial
\[
 \mathcal H_I(u)=P_I(Z)\prod_{v\in I}u_v
           =\sum_{S\subset I}c_S\prod_{v\in S}u_v.
\]
The coefficient \(c_S\) is the expected determinant of the negative
outcome matrix restricted to \(S\).  If \(S\) splits into graph
components, the restricted outcome matrix splits into the same blocks,
and their choices are independent.  Hence \(c_S\) factors over the
components.  Since \(\mathcal H_I(0)=1\), its formal logarithm is
defined.  A coefficient in this logarithm depends only on its supported
subset.  On a disconnected subset \(\mathcal H_I\) factors, so its
logarithm is a sum over components.  Thus every nonconstant monomial
in \(\log\mathcal H_I\) has connected support.

By Equation~\eqref{eq:poly-differential},
\[
                   G_v(Z)=\partial_{z_v}\log P_I(Z).
\]
Apart from the term \(1/z_v\), its monomials arise from connected
supports containing \(v\).  After uniform specialization, the
coefficient of order \(z^{-1-k}\) therefore uses only the radius-\(k\)
ball at \(v\).  The Cauchy-transform expansion identifies that coefficient
with the \(k\)-th rooted moment.  Formal inversion of the series with
leading term \(1/z\) gives the same locality statement for \(F_v\).
Conditioning or changing the shift at \(v\) preserves this argument.

For singleton and two-point sets the coefficients are
\[
 c_{\{v\}}=a_v,\qquad
 c_{\{u,v\}}=a_ua_v-W_uW_v|T_{uv}|^2
                         \mathbb P(\omega_u=\omega_v).
\]
Expanding the logarithmic derivative to order \(z^{-3}\) gives the
first identity in Equation~\eqref{eq:poly-second-moment}.
Alternatively, the arrow matrix has first diagonal \(-a_v\) and
off-diagonal squared norm \(\sum_\lambda b_{v,\lambda}\); comparison
with its second rooted moment gives the second identity.

Under the neighboring-ratio hypothesis, the sum in
Equation~\eqref{eq:poly-second-moment} is at most
\[
 K W_v^2\sum_{u\ne v}|T_{vu}|^2\leq K W_v^2.
\]
This holds for each conditioned law as well, since every collision
probability is at most one.  Also \(a_v\leq\beta tW_v\).
Cauchy--Schwarz for the probability measure \(\rho_v\), followed by
Lemma~\ref{lem:positive-cost}, proves
Equation~\eqref{eq:poly-weight-bounds}.

Finally every outcome has norm at most \((1+\beta t)H_0\), because
principal compression and pinching preserve the contraction bound of
\(T\).  For real \(z>L\), all its characteristic determinants are
positive; for \(z<-L\), all have sign \((-1)^{|I|}\).
Their average therefore has no root outside \([-L,L]\).
Real-rootedness proves Equation~\eqref{eq:poly-root-support}.
The same reasoning applies after restriction or a site option, and
the spectral representation gives the assertion for rooted measures.
\end{proof}

\subsection{The single-site assignment estimate}

We now compare the cost of leaving a site unassigned with the costs of
permanently assigning its possible colors.  The comparison depends only
on the arrow matrix and the averaging of its nonnegative residues.  This
is the finite estimate that will govern the coloring procedure.

\begin{lemma}[Assigning one color]\label{lem:assignment}
For a finite zero-diagonal Hermitian contraction \(T\), arbitrary
positive weights, \(r\geq1\), \(t>0\), and a partial coloring as above,
let \(v\) be unassigned.  With \(\beta=80\), there is
\(i\in\{1,\ldots,r\}\) such that
\begin{equation}
                         d_v^{\,i}-d_v\leq24m_v.
 \label{eq:poly-assignment}
\end{equation}
The bound is independent of the finite set and its weights.
\end{lemma}

\begin{proof}
Write \(a=a_v=tW_v>0\), \(m=m_v\), \(d=d_v\), and
\[
                  H=\begin{pmatrix}-a&g^*\\g&\Lambda\end{pmatrix}
\]
for the old arrow matrix.  By Equation~\eqref{eq:poly-options},
the color-conditioned residues \(b_{i,\lambda}\) are nonnegative and
have average \(b_\lambda\).  A permanent assignment replaces these
residues by \(b_{i,\lambda}\) and the first diagonal by \(-\beta a\).
All deletion eigenvalues \(\Lambda\) stay fixed.

Separate the deletion coordinates into those with \(\lambda>0\) and
those with \(\lambda\leq0\), and use subscripts \(+\) and \(-\) for these
two blocks.  Write \(H=H_+-H_-\) for its positive and negative parts;
subscripts on \(H_\pm\) below specify compressions to the deletion
blocks, not additional positive or negative parts.  Set
\[
 e_+=\Tr(H_-)_{++},\qquad e_-=\Tr(H_+)_{--}.
\]
Since the first diagonal of \(H_+\) is \(m\), comparison of the diagonal
blocks in \(H=H_+-H_-\) gives
\begin{equation}
                       d=m+e_++e_-,
              \qquad (H_-)_{00}=m+a.
 \label{eq:poly-error-decomposition}
\end{equation}
We treat separately the cases in which the rooted cost is at least
the shift and in which it is smaller.

\smallskip\noindent
\emph{Case \(m\geq a\).}
Consider the two side arrows obtained by retaining the first coordinate
and just one deletion block, initially with first diagonal \(-a\).
For the positive block there is exactly one negative eigenvalue,
say \(-L_0\), and
\[
 L_0=a+\sum_{\lambda>0}\frac{b_\lambda}{L_0+\lambda},
             \qquad d^+=L_0-a.
\]
Here \(d^+\) is its positive-trace increment over its deletion block.
The formulas remain valid for an empty positive block, with \(L_0=a\).
For the nonpositive block there is at most one positive eigenvalue.
Let \(s_0\) denote that eigenvalue if present, and put \(s_0=0\)
otherwise.  Its increment is \(d^-=s_0\), and
\begin{equation}
             \sum_{\lambda\leq0}\frac{b_\lambda}{s_0+|\lambda|}
                              \leq s_0+a.
 \label{eq:poly-nonpositive-test}
\end{equation}
If \(s_0=0\), all residues at \(\lambda=0\) must vanish; otherwise the
side quotient would have a positive zero.  Omit those terms in
Equation~\eqref{eq:poly-nonpositive-test}.  If the remaining sum at zero
were greater than \(a\), continuity and the behavior at infinity would
again produce a positive zero, which proves the displayed inequality
also in this case.

Pinching off the other coupling in \(H\) cannot increase its positive
trace.  The uncoupled block contributes exactly its own positive trace,
so \(d^+\leq d\) and \(d^-\leq d\).
By averaging the residues, at least three quarters of the colors satisfy
\[
          \sum_{\lambda>0}\frac{b_{i,\lambda}}{L_0+\lambda}
                                   \leq4d^+,
\]
and at least three quarters satisfy
\[
          \sum_{\lambda\leq0}\frac{b_{i,\lambda}}{s_0+|\lambda|}
                                   \leq4(s_0+a).
\]
Thus some color satisfies both bounds.  A zero old residue is zero for
every color by nonnegativity, so the omitted zero-denominator terms
remain absent.

For this color, the positive-block arrow with first diagonal \(-a\)
has increment at most \(4d^+\).  Indeed, if its negative eigenvalue
magnitude is at most \(L_0\), the assertion is immediate; otherwise
its root equation bounds the increment by the preceding test sum.
The nonpositive-block arrow with first diagonal \(-4a\) has increment
at most \(4s_0\).  If it had a positive root \(s>4s_0\), its decreasing
residue sum would be at most its value at \(s_0\), hence at most
\(4(s_0+a)\), whereas its root equation would require that sum to be
\(s+4a>4(s_0+a)\).

Extend the two side arrows by zero and add them.  The resulting full
arrow has first diagonal \(-5a\).  Subadditivity of positive trace,
which follows from its variational formula, shows that its increment
is at most \(4d^++4d^-\).  Decreasing the first diagonal further to
\(-\beta a\) can only decrease positive trace.  Therefore
\[
       d_v^{\,i}\leq4(d^++d^-)\leq8d,\qquad
       d_v^{\,i}-d\leq7d\leq7(2m+a)\leq21m,
\]
where Lemma~\ref{lem:positive-cost} was used in the last step.

\smallskip\noindent
\emph{Case \(m<a\).}
The positive and negative parts of \(H\) give more useful residue
bounds in this case.  For a positive block matrix, the outer product
of its first coupling column is bounded by its first diagonal entry
times its remaining block.  Apply this to \(H_+\) and \(H_-\), and use
\((p-q)(p-q)^*\leq2pp^*+2qq^*\) for the difference of their coupling
columns.  As
\((H_+)_{++}=\Lambda_++(H_-)_{++}\) and
\((H_-)_{--}=|\Lambda_-|+(H_+)_{--}\), this gives
\begin{align}
 g_+g_+^*&\leq
       2m\Lambda_++(4m+2a)(H_-)_{++},
 \label{eq:poly-positive-block}\\
 g_-g_-^*&\leq
       2(m+a)|\Lambda_-|+(4m+2a)(H_+)_{--}.
 \label{eq:poly-negative-block}
\end{align}
Conjugate Equation~\eqref{eq:poly-positive-block} by
\((32a+\Lambda_+)^{-1/2}\) and take the operator norm.
The left side has rank at most one, so its norm is the corresponding
residue sum.  The first term on the right has norm at most \(2m\),
and the error term has norm at most its trace.  Since \(4m+2a\leq6a\),
\begin{equation}
                 \sum_{\lambda>0}\frac{b_\lambda}{32a+\lambda}
                              \leq2m+\frac6{32}e_+.
 \label{eq:poly-positive-test}
\end{equation}

Put \(L_1=e_-\).  If \(L_1>0\), conjugating
Equation~\eqref{eq:poly-negative-block} by
\((L_1+|\Lambda_-|)^{-1/2}\) in the same way gives
\begin{equation}
 \sum_{\lambda\leq0}\frac{b_\lambda}{L_1+|\lambda|}
              \leq2(m+a)+(4m+2a)\leq10a.
 \label{eq:poly-negative-test}
\end{equation}
Indeed the trace of \((H_+)_{--}\) is \(L_1\), so its conjugated norm
is at most \(1\).  If \(L_1=0\), positivity implies
\((H_+)_{--}=0\).  Equation~\eqref{eq:poly-negative-block} then forces
zero coupling on the kernel of \(|\Lambda_-|\).  Use its inverse only
on its nonzero support and omit the zero residues at \(\lambda=0\);
Equation~\eqref{eq:poly-negative-test} remains valid.

Choose a color for which both residue sums in
Equations~\eqref{eq:poly-positive-test}
and~\eqref{eq:poly-negative-test} are at most four times their indicated
upper bounds.  Such a color exists by the same averaging argument.
For that color, the positive-block arrow with diagonal \(-32a\) has
negative eigenvalue magnitude at least \(32a\); its root equation
therefore bounds its positive-trace increment by
\[
                              8m+\frac{24}{32}e_+.
\]
The nonpositive-block arrow with diagonal \(-40a\) has increment at
most \(L_1\).  Indeed a positive root \(s>L_1\) would have residue sum
at most its test value \(40a\), whereas its root equation requires
\(s+40a>40a\).  This also covers \(L_1=0\), with zero residues omitted
as above.

Adding the side arrows gives the full arrow with diagonal \(-72a\).
Subadditivity of positive trace and then decreasing the diagonal to
\(-80a\) yield
\[
                 d_v^{\,i}\leq8m+\frac34e_++e_-.
\]
Together with Equation~\eqref{eq:poly-error-decomposition}, this gives
\[
                      d_v^{\,i}-d\leq7m-\frac14e_+\leq7m.
\]
Both cases imply Equation~\eqref{eq:poly-assignment}.  All arguments
allow empty deletion blocks and empty sums, so they include a
one-site matrix and every zero-residue case.
\end{proof}
